\documentclass[10pt,a4paper]{amsart}
\usepackage[margin=19mm]{geometry}
\usepackage{amsmath,amssymb,mathtools}
\usepackage[scr=rsfs,cal=euler]{mathalfa}
\usepackage{xfrac}
\usepackage[unicode,hidelinks,bookmarks=false]{hyperref}
\newcommand{\ie}{i.e.\ }
\newcommand{\cf}{cf.\ }
\newcommand{\resp}{respectively\ }
\newcommand{\Subsection}{\subsection}
\providecommand{\nobreakdash}{\nobreak\hbox{-}}
\setlength{\emergencystretch}{2em}
\multlinegap0pt
\usepackage{bm}
\usepackage{bbm}
\usepackage{dsfont}
\usepackage{xspace}
\usepackage{cancel}
\usepackage{mathtools}
\usepackage{amsthm}
\makeatletter
\@ifundefined{Def}{\newtheorem{Def}{Def}}{}
\@ifundefined{lem}{\newtheorem{lem}[Def]{Lemma}}{}
\@ifundefined{thm}{\newtheorem{thm}[Def]{Theorem}}{}
\makeatother
\newcommand{\C}{{\mathbb C}}
\newcommand{\Q}{{\mathbb Q}}
\newcommand{\Rea}{\operatorname{Re}}
\newcommand{\R}{{\mathbb R}}
\newcommand{\SL}{\operatorname{{\mathbf SL}}}
\newcommand{\Z}{{\mathbb Z}}
\title[Periodic Hypersurfaces and Lee-Yang Polynomials]{Periodic Hypersurfaces and Lee-Yang Polynomials}
\author{Lior Alon, Mario Kummer}
\date{Source: arXiv:2507.16029v3 (CC BY 4.0)}
\begin{document}
\maketitle

\noindent\emph{Source and licence.} Periodic Hypersurfaces and Lee-Yang Polynomials. Version arXiv:2507.16029v3 (CC BY 4.0), distributed under the Creative Commons Attribution 4.0 International licence, as recorded in the arXiv licence metadata for this exact version and shown on the abstract page https://arxiv.org/abs/2507.16029v3. Full rights notice: licenses/arxiv-2507.16029-CC-BY-4.0.txt.\par\smallskip

\noindent\textit{Editorial scope.} Counted unit: the Thm whose source label is \texttt{thm:essleeyangchara} in the article (source0.tex lines 772--780, source label \texttt{thm:essleeyangchara}), delivered together with the article's own complete original proof of it (source0.tex lines 781--798, one \texttt{proof} environment of 18 lines). No mathematics is written, repaired or re-derived: statement and proof are the article's own text line for line. Required context retained uncounted: lem:denselinerealrootedessly (lines 760--762). Every definition, display and label of those lines is retained unchanged. Omitted from this package and not redistributed: 1--759, 763--771, 799--1044. Nothing inside a retained excerpt is omitted or altered: every retained body is delivered exactly as the article's lines are written, so no omission receipt of any kind accompanies this package. Citations: the retained excerpts carry no citation command at all, so no bibliography entry of the article is transcribed and no reference is redistributed. Reference adaptation: none was needed and none was made; every reference inside the delivered unit resolves inside it, so no unresolved reference or citation is produced. Printed slips kept literal: no printed word, symbol, hypothesis or number of the counted statement or of its proof was altered. Layout and class: the article's own class is \texttt{amsart} with packages including amsrefs, amsfonts, mathrsfs, amsthm, amssymb, enumerate, tikz-cd, cleveref, enumitem; the wrapper is the lane's pinned \texttt{amsart} editorial wrapper at 10pt on A4, which restates the theorem-like environments the retained text needs and adds the article's own macro definitions that the retained text needs (\texttt{\textbackslash C}, \texttt{\textbackslash Q}, \texttt{\textbackslash R}, \texttt{\textbackslash Rea}, \texttt{\textbackslash SL}, \texttt{\textbackslash Z}), with extra wrapper packages (bm, bbm, dsfont, xspace, cancel, mathtools). The delivered numbering is the unit's own. Assets: no retained span contains a figure environment, a graphics-inclusion command, a PDF-page inclusion, a table or a TikZ picture; no image file is copied into this package, the package declares no asset dependency and no image file is redistributed with it. Licence: version arXiv:2507.16029v3 of this article is distributed under the Creative Commons Attribution 4.0 International licence, as recorded in the arXiv licence metadata for this exact version and shown on the abstract page https://arxiv.org/abs/2507.16029v3. A full rights notice is delivered at licenses/arxiv-2507.16029-CC-BY-4.0.txt. Characters: every retained excerpt is pure ASCII, so the delivered engine, XeLaTeX with its default Latin Modern text font, reports no missing character.
\begin{lem}\label{lem:denselinerealrootedessly}
    Let $p\in\C[z_1^{\pm},\ldots,z_n^{\pm}]$ be a Laurent polynomial and $\ell\in\R^n$ a vector with $\Q$-linearly independent entries such that $p(\exp(2\pi i\ell t))$ has only real zeros. There exists $k\in\Z^n$, a matrix $A\in\SL_n(\Z)$, a Lee--Yang polynomial $q\in\C[z_1,\ldots,z_n]$ and a vector $\tilde\ell\in\R^n$ with positive entries such that $q(z)=z^k\cdot p(z^A)$ and $\ell=A\tilde\ell$.
\end{lem}

\begin{thm}\label{thm:essleeyangchara}
    Let $p\in\C[z_1^{\pm},\ldots,z_n^{\pm}]$ be a Laurent polynomial. The following are equivalent:
    \begin{enumerate}
        \item $p$ is essentially Lee--Yang.
        \item There exists an open cone $\emptyset\neq K\subset\R^n$ such that $p(\exp(2\pi i\ell t)$ is real-rooted for all $\ell\in K$.
        \item There exists a vector $\ell\in\R^n$ with $\Q$-linearly independent entries such that $p(\exp(2\pi i\ell t))$ has only real zeros.
        \item There exists $k\in\Z^n$, a matrix $A\in\SL_n(\Z)$ and a Lee--Yang polynomial $q\in\C[z_1,\ldots,z_n]$ with $q(z)=z^kp(z^A)$.
    \end{enumerate}
\end{thm}

\begin{proof}
    ``$(1)\Rightarrow(2)$'': By definition there exists a nonempty open cone $K\subset\R^n$ such that $p(\exp(z))\neq0$ for all $z\in\C^n$ with $\Rea(z)\in -K\cup K$. Let $\ell\in K$ and $t_0\in\C$ with non-zero imaginary part $0\neq b\in\R$. Then $\Rea(2\pi i\ell t_0)=-2b\pi \ell\in -K\cup K$. Hence $t$ is not a zero of $p(\exp(2\pi i\ell t)$ which is thus real-rooted.

    ``$(2)\Rightarrow(3)$'': This is clear because every non-empty open set contains a point with $\Q$-linearly independent entries.

    ``$(3)\Rightarrow(4)$'': This is part of Lemma \ref{lem:denselinerealrootedessly}.

    ``$(4)\Rightarrow(1)$'':  Consider the open cone $K=A\R_{+
    }^n$. Let $z\in\C^n$ with $\Rea(z)\in A\R_{+
    }^n\cup A\R_{-
    }^n$. Letting $z'=A^{-1}z\in\C^n$, we have $\Rea(z')=A^{-1}\Rea(z)\in \R_{+
    }^n\cup \R_{-
    }^n$ and
    \begin{equation*}
        p(\exp(z))=p(\exp(Az'))=e^{-\langle  k,z'\rangle}q(\exp(z'))\neq0
    \end{equation*}
    because $q$ is Lee--Yang.
\end{proof}

\end{document}
